% Versioned self-contained section source; integrated verbatim into fixed124.
\subsection{Dyadic Fox--Bockstein--orientation completion}
\label{subsec:dyadic-fox-orientation}

The odd-prime symplectic normal form does not extend verbatim to
$p=2$: a square term already contributes in Magnus degree two and
need not vanish in the cup form.  More importantly, for the
exceptional Demu\v{s}kin invariant $q=2$, the quadratic Fox form and
the ordinary integral exponent-sum/Bockstein information do not
record the canonical dualizing-orientation image.  We now close this
dyadic classification gap without changing the finite Fox
resolution or the all-Sylow descent \cite{LabuteClassification,NSW}.

\begin{theorem}[Four dyadic Demu\v{s}kin forms and their quadratic symbols]
\label{thm:dyadic-relator-cases}
Let $P$ be an infinite finitely generated Demu\v{s}kin pro-$2$
group of rank $d$, with invariant $q=q(P)$, and write
$[x,y]=x^{-1}y^{-1}xy$.  In a suitable minimal presentation
$P=\langle x_1,\ldots,x_d\mid r\rangle_{\mathrm{pro}-2}$,
the relation $r$ is in exactly one of the following normal-form
families:
\begin{enumerate}[label=\textnormal{(\roman*)}]
\item If $q\ne2$, then $d$ is even and
\[
r=x_1^q[x_1,x_2][x_3,x_4]\cdots[x_{d-1},x_d],
\qquad q=0\ \text{or}\ q=2^f,\ f\ge2.
\]
\item If $q=2$, $d$ is odd and $d\ge3$, then
\[
r=x_1^2x_2^{2^f}[x_2,x_3][x_4,x_5]\cdots[x_{d-1},x_d],
\quad f\in\{2,3,\ldots,\infty\}.
\]
\item If $q=2$, $d$ is even and the canonical-orientation
image has square-index $2$, then
\[
r=x_1^{2+2^f}[x_1,x_2][x_3,x_4]\cdots[x_{d-1},x_d],
\quad f\in\{2,3,\ldots,\infty\}.
\]
\item If $q=2$, $d\ge4$ is even and the canonical-orientation
image has square-index $4$, then
\[
r=x_1^2[x_1,x_2]x_3^{2^f}[x_3,x_4]\cdots[x_{d-1},x_d],
\quad f\in\{2,3,\ldots\}.
\]
\end{enumerate}
Here $2^\infty=0$ means that the corresponding factor is absent.
With $C_{ij}=X_iX_j+X_jX_i\in\mathbf F_2\langle X_1,\ldots,X_d\rangle$,
their degree-two Magnus symbols are
\begin{equation}\label{eq:dyadic-quadratic-symbols}
\sigma_2(r)=
\begin{cases}
\displaystyle\sum_{j=1}^{d/2}C_{2j-1,2j}, &(i),\\[2pt]
\displaystyle X_1^2+\sum_{j=1}^{(d-1)/2}C_{2j,2j+1},
 &(ii),\\[2pt]
\displaystyle X_1^2+\sum_{j=1}^{d/2}C_{2j-1,2j},
 &(iii),(iv).
\end{cases}
\end{equation}
The induced Fox cup form is nondegenerate, but in cases
(ii)--(iv) it is \emph{nonalternating}: $\chi_1^2=\omega$.

For $P=G_K(2)$ with $K/\mathbf Q_2$ finite,
one always has $q\ge2$, $d=[K:\mathbf Q_2]+2$,
and an open orientation image, so any occurring $f$ is finite.
\end{theorem}

\begin{proof}
The four presentations and the square-index distinction are
the Demu\v{s}kin--Serre--Labute classification
\cite{LabuteClassification,NSW}.
Over $\mathbf F_2$ one has
$\mathcal M([x_i,x_j])=1+C_{ij}+O(\deg3)$ and
$\mathcal M(x_i^{2^f})=1+X_i^{2^f}$.
For $f\ge2$ the latter does not enter quadratic degree.
The factors $x_1^2$ and $x_1^{2+2^f}$, however,
both contribute $X_1^2$.
This proves \eqref{eq:dyadic-quadratic-symbols}.
The quadratic Fox cup theorem
(Theorem~\ref{thm:quadratic-fox-cup})
identifies its coefficients with the cup matrix.
Its blocks are symplectic $2$-blocks together,
in the exceptional cases, with either the isolated
diagonal $1$ or the nonalternating block
$\bigl(\begin{smallmatrix}1&1\\1&0\end{smallmatrix}\bigr)$;
in all cases the form is nondegenerate.
For the local assertion, $\mu_2\subset K$ implies that
$G_K(2)$ is Demu\v{s}kin, and its dualizing character
is cyclotomic.  Its image in $\mathbb Z_2^\times$ is open
because $K$ is finite.
\end{proof}

Let $\partial_i$ be the completed \emph{left} Fox derivative,
and for a continuous character $\theta:P\to\mathbb Z_2^\times$
write
\[
\varepsilon_\theta:\mathbb Z_2[[F_d]]\longrightarrow\mathbb Z_2,
\quad x_i\longmapsto\theta(x_i)
\]
for the associated twisted augmentation.  The system
\begin{equation}\label{eq:dyadic-twisted-fox-orientation}
\boxed{\quad\varepsilon_\theta(\partial_i r)=0
\quad(1\le i\le d)\quad}
\end{equation}
says that every continuous $\theta$-crossed derivation of
the free pro-$2$ group vanishes on the relator.  It is
presentation-covariant, although its coordinates are not
generator-independent.

\begin{theorem}[Canonical dyadic orientation from twisted Fox equations]
\label{thm:dyadic-twisted-fox-orientation}
On each normal form of Theorem~\ref{thm:dyadic-relator-cases},
Equation~\eqref{eq:dyadic-twisted-fox-orientation} admits
a unique solution in generator values
$\theta(x_i)\in\mathbb Z_2^\times$.
It is the canonical dualizing orientation $\theta_P$.
All unmentioned generator values are $1$, and the nontrivial
values are
\begin{align}
\textnormal{(i)}\quad&
 \theta(x_1)=1,\quad
 \theta(x_2)=(1-q)^{-1};\label{eq:dyadic-orientation-i}\\
\textnormal{(ii)}\quad&
 \theta(x_1)=-1,\quad\theta(x_2)=1,\quad
 \theta(x_3)=(1-2^f)^{-1};\label{eq:dyadic-orientation-ii}\\
\textnormal{(iii)}\quad&
 \theta(x_1)=1,\quad
 \theta(x_2)=(-1-2^f)^{-1};
 \label{eq:dyadic-orientation-iii}\\
\textnormal{(iv)}\quad&
 \theta(x_1)=1,\quad\theta(x_2)=-1,\quad
 \theta(x_3)=1,\quad\theta(x_4)=(1-2^f)^{-1}.
 \label{eq:dyadic-orientation-iv}
\end{align}
The corresponding closed orientation images are
\begin{equation}\label{eq:dyadic-orientation-images}
\begin{aligned}
\textnormal{(i)}\quad&
 1+q\mathbb Z_2\ (q\ge4),\quad\{1\}\ (q=0),\\
\textnormal{(ii)}\quad&
 \{\pm1\}(1+2^f\mathbb Z_2),\\
\textnormal{(iii)}\quad&
 \overline{\langle(-1-2^f)^{-1}\rangle},\\
\textnormal{(iv)}\quad&
 \{\pm1\}(1+2^f\mathbb Z_2).
\end{aligned}
\end{equation}
For $f=\infty$ use $1+2^f\mathbb Z_2=\{1\}$ and
$(-1-2^f)^{-1}=-1$.
\end{theorem}

\begin{proof}
With the commutator convention above,
\[
\varepsilon_\theta(\partial_x[x,y])
=\theta(x)^{-1}(\theta(y)^{-1}-1),
\quad
\varepsilon_\theta(\partial_y[x,y])
=\theta(x)^{-1}\theta(y)^{-1}(\theta(x)-1).
\]
Multiplication by a preceding word only adds a unit factor.
A pure commutator block therefore forces both orientation
values to $1$.  In a block $x^m[x,y]$, the $y$-derivative
forces $\theta(x)=1$, after which the $x$-derivative
becomes $m+\theta(y)^{-1}-1$; hence
$\theta(y)=(1-m)^{-1}$.
For the odd-rank family, the first derivative of $x_1^2$
forces $\theta(x_1)=-1$, then the $(x_2,x_3)$ block
forces $\theta(x_2)=1$ and
$\theta(x_3)=(1-2^f)^{-1}$.
For family (iv), the first squared commutator block gives
$(\theta(x_1),\theta(x_2))=(1,-1)$,
followed by the ordinary $x_3^{2^f}[x_3,x_4]$ block.
This gives the four displayed lists and proves uniqueness.
The Fox identity $r-1=\sum_i\partial_i r\,(x_i-1)$
ensures that the resulting map kills $r$, and
the known canonical orientation in Labute's classification
has precisely these normal-form values
\cite{LabuteClassification,NSW}.
For finite $f$, $(1-2^f)^{-1}$ generates
$1+2^f\mathbb Z_2$ topologically, as does
$(1-q)^{-1}$ when $q\ge4$.
The remaining image claims are immediate.
\end{proof}

\begin{corollary}[Finite classification-faithful dyadic symbol]
\label{thm:dyadic-classification-certificate}
The presentation-covariant certificate
\begin{equation}\label{eq:dyadic-classification-certificate}
\boxed{\quad
\mathfrak C_2(P)=
\bigl(d,\,[\sigma_2(r)],\,
\mathfrak q(P),\,\operatorname{im}\theta_P\bigr)
\quad}
\end{equation}
determines the isomorphism class of every
Demu\v{s}kin pro-$2$ group in
Theorem~\ref{thm:dyadic-relator-cases}.
Here $\mathfrak q(P)$ is the ideal of the augmented
integral first Fox derivatives and $\theta_P$
is determined by the twisted Fox system.
On all $q=2$ normal forms one has
$\mathfrak q(P)=(2)$ and the same diagonal-square
information; the orientation image records the extra
Labute index and parameter $f$.
\end{corollary}

\begin{proof}
The quadratic Fox theorem gives the mod-$2$ cup form,
and the augmented first derivatives give the
abelianization's torsion exponent ideal.
Theorem~\ref{thm:dyadic-twisted-fox-orientation}
recovers the canonical orientation.
Labute's classification makes $(d,\operatorname{im}\theta_P)$
a complete isomorphism invariant in the finite-rank
Demu\v{s}kin sector \cite{LabuteClassification};
the displayed richer certificate is therefore faithful.
The individual matrix coordinates need not be
presentation-invariant.
\end{proof}

\begin{proposition}[The explicit $G_{\mathbf Q_2}(2)$ anchor]
\label{prop:dyadic-q2-explicit-anchor}
Let
\[
P_{\mathbf Q_2}=G_{\mathbf Q_2}(2)
=\langle a,s,y\mid a^2s^4[s,y]\rangle_{\mathrm{pro}-2}.
\]
For the Magnus variables $A,S,Y$, the mod-$2$
Fox--Cartan and integral Fox--Bockstein data are
\begin{equation}\label{eq:dyadic-q2-anchor-data}
\begin{gathered}
\sigma_2(r)=A^2+SY+YS,\qquad
B_{\rm cup}=
 \begin{pmatrix}1&0&0\\0&0&1\\0&1&0\end{pmatrix},\\
\lambda(r)=(2,4,0),\quad \mathfrak q(P_{\mathbf Q_2})=(2),\\
\beta_1(\chi_a)=\omega,\qquad
\beta_1(\chi_s)=\beta_1(\chi_y)=0,\\
(\theta(a),\theta(s),\theta(y))=
(-1,1,-1/3),\quad
\operatorname{im}\theta=\mathbb Z_2^\times.
\end{gathered}
\end{equation}
For trivial integral coefficients, its finite
Fox cochain complex is
\[
\mathbb Z_2\xrightarrow{0}\mathbb Z_2^3
 \xrightarrow{(2,4,0)}\mathbb Z_2.
\]
Its lifting filtration satisfies
$L_1=H^1(P_{\mathbf Q_2},\mathbf F_2)$ and
$L_m=\operatorname{span}_{\mathbf F_2}
\{\chi_s,\chi_y\}$ for every $m\ge2$.
\end{proposition}

\begin{proof}
The displayed relator is the rank-$3$, $q=2$,
$f=2$ normal form for $G_{\mathbf Q_2}(2)$
\cite{LabuteClassification,NSW,RoeTurtureanGQ2}.
Modulo $2$, $a^2$ contributes $A^2$,
$s^4$ has no quadratic term, and $[s,y]$
contributes $SY+YS$.
The quadratic Fox cup theorem gives the matrix.
The exponent sums are $(2,4,0)$.
In the integral trivial-coefficient Fox complex,
a $1$-cochain $(u,v,w)$ is closed modulo $2^m$
exactly when $2u+4v\equiv0\pmod{2^m}$.
For $m=2$, dividing its coboundary by $2$
gives the ordinary Bockstein $(1,0,0)$ mod $2$;
for all $m\ge2$, precisely the reductions with
$u=0$ lift, proving the filtration.
Finally the twisted Fox equations imply
$\theta(a)=-1$, $\theta(s)=1$ and
$4+\theta(y)^{-1}-1=0$, so
$\theta(y)=-1/3$.
Its image together with $-1$ is
$\{\pm1\}(1+4\mathbb Z_2)=\mathbb Z_2^\times$.
\end{proof}

\begin{warning}[The ordinary dyadic Bockstein depth is insufficient]
\label{warn:dyadic-bockstein-depth-insufficient}
For each finite $f\ge2$, the rank-three relations
\[
r_f=x_1^2x_2^{2^f}[x_2,x_3]
\]
have the same $\sigma_2(r_f)$, integral
Fox ideal $(2)$, first Bockstein, and full
untwisted lifting filtration
\[
L_1=H^1(P_f,\mathbf F_2),\qquad
L_m=\operatorname{span}\{\chi_2,\chi_3\}
\quad(m\ge2).
\]
Nevertheless their orientation images
$\{\pm1\}(1+2^f\mathbb Z_2)$ vary with $f$.
Thus the original odd-prime symbol
$(d,\sigma_2,\mathfrak q)$ does not classify
the exceptional dyadic groups.
The missing input is the \emph{twisted}
Fox orientation, not another untwisted
Bockstein depth.
\end{warning}

\begin{remark}[Interface with all-Sylow finite closure]
\label{rem:dyadic-all-sylow-interface}
For any residual sector at $p=2$, the preimage
$H$ of a residual Sylow $2$-subgroup has odd
index in $G_K$ and is the absolute Galois
group of a finite dyadic extension $K_H$.
Because $\mu_2\subset K_H$, the maximal
pro-$2$ quotient $H(2)$ is always
Demu\v{s}kin, not the free branch.
Choosing the normal form determined by its
cyclotomic orientation and evaluating its
completed Fox derivatives gives exactly
the bounded finite-free coefficient chart
of Theorem~\ref{thm:family-sylow-fox}.
The odd index supplies the integral
transfer idempotent from
Theorem~\ref{thm:sylow-fox-retract}.
Hence the additional dyadic classification
data alter only the explicit finite
Fox matrices, not perfectness, derived base
change, Cayley--Hamilton transport, or
the recursive globalization theorem.
The maximal pro-$2$ presentation is not
misidentified with a marked presentation
of the entire absolute Galois group.
\end{remark}
